\documentclass[a4paper,11pt]{ltjsarticle}
\usepackage{luatexja}
\usepackage{amsmath,amssymb,amsthm}
\usepackage{tikz}
\usepackage{tikz-cd}
\usetikzlibrary{positioning,arrows.meta,shapes.geometric}
\usepackage{tcolorbox}
\tcbuselibrary{theorems,skins,breakable}
\usepackage{hyperref}
\usepackage{listings}
\usepackage{xcolor}

% カラースキーム
\definecolor{maincolor}{RGB}{70,130,180}
\definecolor{examplecolor}{RGB}{60,179,113}
\definecolor{codebackground}{RGB}{245,245,245}

% 定理環境
\newtheorem{theorem}{定理}[section]
\newtheorem{definition}[theorem]{定義}
\newtheorem{proposition}[theorem]{命題}
\newtheorem{lemma}[theorem]{補題}
\newtheorem{corollary}[theorem]{系}

\newtcolorbox{examplebox}[1][]{
  colback=examplecolor!5!white,
  colframe=examplecolor!75!black,
  fonttitle=\bfseries,
  title=例,
  #1
}

\newtcolorbox{remarkbox}[1][]{
  colback=maincolor!5!white,
  colframe=maincolor!75!black,
  fonttitle=\bfseries,
  title=注意,
  #1
}

% Leanコードスタイル
\lstdefinestyle{leanstyle}{
  language=ML,
  backgroundcolor=\color{codebackground},
  basicstyle=\ttfamily\small,
  keywordstyle=\color{blue}\bfseries,
  commentstyle=\color{gray}\itshape,
  stringstyle=\color{red},
  numbers=left,
  numberstyle=\tiny\color{gray},
  stepnumber=1,
  numbersep=5pt,
  frame=single,
  breaklines=true,
  breakatwhitespace=false,
  tabsize=2,
  captionpos=b
}

\title{\Huge \textbf{構造塔理論入門} \\[0.5em]
\Large Bourbaki的母構造の現代的形式化}

\author{
\textbf{形式化：} Claude Code (Anthropic) によるTeXファイル生成 \\[0.3em]
\textbf{実装：} CODEX (OpenAI) によるLean 4形式化 \\[0.3em]
\textbf{プロジェクト：} Lean Natural Numbers - Structure Tower Formalization
}

\date{\today}

\begin{document}

\maketitle

\begin{abstract}
本稿は、Nicolas Bourbakiが提唱した数学の構造主義的アプローチを、Lean 4定理証明支援系を用いて形式化した「構造塔理論」を学部生向けに解説する。構造塔は、数学的対象の階層的な包含関係を統一的に扱うための枠組みであり、位相構造、代数構造、順序構造といったBourbaki的母構造のすべてに適用可能である。本稿では、整数の絶対値階層、有限集合の濃度、部分群の包含関係など、具体的な数学的構造を通じて構造塔の概念を説明し、抽象理論と具体例の架け橋を提供する。
\end{abstract}

\tableofcontents
\newpage

\section{導入：なぜ構造塔なのか}

\subsection{数学における「構造」の統一的理解}

20世紀の数学者集団Nicolas Bourbakiは、数学を「構造の研究」として捉え直すことを提唱した。彼らは、一見異なる数学的対象（群、位相空間、順序集合など）が、実は共通の「構造」を持っていることを認識し、数学を以下の3つの**母構造**（structures mères）に分類した：

\begin{enumerate}
\item \textbf{代数構造}：群、環、体など（演算の性質）
\item \textbf{順序構造}：順序集合、束など（順序関係）
\item \textbf{位相構造}：位相空間、連続性など（近さの概念）
\end{enumerate}

しかし、これらの構造を\textbf{階層的に扱う}統一的な枠組みは、Bourbakiの時代には明示的に形式化されていなかった。

\subsection{構造塔：階層性の形式化}

\textbf{構造塔}（Structure Tower）は、数学的対象の**包含関係の階層**を抽象化した概念である。例えば：

\begin{itemize}
\item 整数 $\mathbb{Z}$ において、$\{k \in \mathbb{Z} : |k| \leq 3\} \subseteq \{k \in \mathbb{Z} : |k| \leq 5\}$
\item 群 $G$ において、部分群 $H_1 \subseteq H_2 \subseteq G$
\item 位相空間 $X$ において、粗い位相 $\subseteq$ 細かい位相
\end{itemize}

これらはすべて、「層」（layer）と呼ばれる部分集合が、ある添字集合上の順序に従って単調に増大する構造を持つ。構造塔は、この共通性を捉える。

\begin{remarkbox}
構造塔の特徴は、各数学的対象に対して「最小の層」が存在することである。これにより、対象を適切に分類し、階層的な理解が可能になる。
\end{remarkbox}

\newpage

\section{構造塔の数学的定義}

\subsection{形式的定義}

\begin{definition}[構造塔]
集合 $X$ と半順序集合 $(I, \leq)$ に対して、\textbf{最小層を持つ構造塔}（Structure Tower with Minimal Layer）は以下のデータからなる：

\begin{enumerate}
\item \textbf{層関数} $\text{layer} : I \to \mathcal{P}(X)$（$\mathcal{P}(X)$ は $X$ の冪集合）
\item \textbf{被覆性}：$\forall x \in X, \exists i \in I$ such that $x \in \text{layer}(i)$
\item \textbf{単調性}：$\forall i, j \in I, \ i \leq j \Rightarrow \text{layer}(i) \subseteq \text{layer}(j)$
\item \textbf{最小層関数} $\text{minLayer} : X \to I$
\item \textbf{最小層の帰属性}：$\forall x \in X, \ x \in \text{layer}(\text{minLayer}(x))$
\item \textbf{最小性}：$\forall x \in X, \forall i \in I, \ x \in \text{layer}(i) \Rightarrow \text{minLayer}(x) \leq i$
\end{enumerate}
\end{definition}

\subsection{定義の直観的理解}

\begin{center}
\begin{tikzpicture}[
  layer/.style={draw, ellipse, minimum width=3cm, minimum height=1.5cm, fill=maincolor!20},
  point/.style={circle, fill, inner sep=2pt}
]
% 層の描画
\node[layer] (L1) at (0,0) {$\text{layer}(1)$};
\node[layer, minimum width=4.5cm] (L2) at (0,2) {$\text{layer}(2)$};
\node[layer, minimum width=6cm] (L3) at (0,4) {$\text{layer}(3)$};

% 点の配置
\node[point, label=below:$x_1$] at (-1,0.2) {};
\node[point, label=below:$x_2$] at (0.5,0.2) {};
\node[point, label=right:$x_3$] at (1.5,2.2) {};
\node[point, label=right:$x_4$] at (-1.8,4.2) {};

% 注釈
\draw[->,thick] (L1) -- (L2) node[midway,right] {単調性};
\draw[->,thick] (L2) -- (L3);

\node[right=4cm of L2, text width=4cm, align=left] {
  各点 $x$ は最小の層 $\text{layer}(\text{minLayer}(x))$ に属する。\\[0.5em]
  例：$x_1 \in \text{layer}(1)$ なら\\
  $\text{minLayer}(x_1) = 1$
};
\end{tikzpicture}
\end{center}

\begin{remarkbox}[title=重要な性質]
\begin{itemize}
\item 被覆性により、すべての要素がどこかの層に属する
\item 単調性により、小さい層に属する要素は大きい層にも属する
\item 最小層により、各要素を一意に分類できる
\end{itemize}
\end{remarkbox}

\newpage

\section{具体例1：整数の絶対値による構造塔}

\subsection{構成}

整数全体 $\mathbb{Z}$ に対して、以下のように構造塔を定義する：

\begin{itemize}
\item $X = \mathbb{Z}$（台集合）
\item $I = \mathbb{N}$（自然数、通常の順序）
\item $\text{layer}(n) = \{k \in \mathbb{Z} : |k| \leq n\}$
\item $\text{minLayer}(k) = |k|$（絶対値）
\end{itemize}

\subsection{視覚化}

\begin{center}
\begin{tikzpicture}[scale=0.8]
% 層0
\draw[fill=maincolor!10] (-0.5,-0.5) rectangle (0.5,0.5);
\node at (0,0) {$0$};
\node[left] at (-0.5,0) {層 0:};

% 層1
\draw[fill=maincolor!20] (-1.5,-1.5-1) rectangle (1.5,0.5-1);
\node at (-1,-1) {$-1$};
\node at (0,-1) {$0$};
\node at (1,-1) {$1$};
\node[left] at (-1.5,-1) {層 1:};

% 層2
\draw[fill=maincolor!30] (-2.5,-1.5-3) rectangle (2.5,0.5-3);
\node at (-2,-3) {$-2$};
\node at (-1,-3) {$-1$};
\node at (0,-3) {$0$};
\node at (1,-3) {$1$};
\node at (2,-3) {$2$};
\node[left] at (-2.5,-3) {層 2:};

% 層3
\draw[fill=maincolor!40] (-3.5,-1.5-5) rectangle (3.5,0.5-5);
\node at (-3,-5) {$-3$};
\node at (-2,-5) {$-2$};
\node at (-1,-5) {$-1$};
\node at (0,-5) {$0$};
\node at (1,-5) {$1$};
\node at (2,-5) {$2$};
\node at (3,-5) {$3$};
\node[left] at (-3.5,-5) {層 3:};

% 包含関係
\draw[->,very thick] (0.6,0) -- (0.6,-1.4) node[midway,right] {$\subseteq$};
\draw[->,very thick] (0.6,-1.6) -- (0.6,-2.4);
\draw[->,very thick] (0.6,-3.6) -- (0.6,-4.4);
\end{tikzpicture}
\end{center}

\subsection{Lean 4 での実装}

\begin{lstlisting}[style=leanstyle, caption=整数絶対値塔の定義]
def IntAbsoluteTower : StructureTowerMin ℤ ℕ where
  layer n := {k : ℤ | k.natAbs ≤ n}
  covering := by
    intro x
    exact ⟨x.natAbs, by simp⟩
  monotone := by
    intro m n hmn k hk
    exact le_trans hk hmn
  minLayer := Int.natAbs
  minLayer_mem := by
    intro x
    simp
  minLayer_minimal := by
    intro x n hn
    exact hn
\end{lstlisting}

\begin{examplebox}
\textbf{具体的な計算：}
\begin{itemize}
\item $\text{minLayer}(6) = 6$ より、$6 \in \text{layer}(6)$
\item $\text{minLayer}(-3) = 3$ より、$-3 \in \text{layer}(3)$
\item 単調性より、$-3 \in \text{layer}(5)$ も成り立つ
\end{itemize}
\end{examplebox}

\newpage

\section{具体例2：有限集合の濃度による構造塔}

\subsection{構成}

型 $\alpha$ 上の有限集合 $\mathsf{Finset}(\alpha)$ に対して：

\begin{itemize}
\item $X = \mathsf{Finset}(\alpha)$
\item $I = \mathbb{N}$
\item $\text{layer}(n) = \{s \in \mathsf{Finset}(\alpha) : |s| \leq n\}$
\item $\text{minLayer}(s) = |s|$（濃度）
\end{itemize}

\subsection{視覚化}

\begin{center}
\begin{tikzpicture}[
  set/.style={draw, circle, minimum size=0.8cm},
  level/.style={rectangle, draw=maincolor, fill=maincolor!10, minimum width=8cm, minimum height=1.5cm}
]
% 層0
\node[level] (L0) at (0,0) {};
\node[left] at (-4,0) {層 0:};
\node[set] at (0,0) {$\emptyset$};

% 層1
\node[level] (L1) at (0,-2) {};
\node[left] at (-4,-2) {層 1:};
\node[set] at (-2,-2) {$\{a\}$};
\node[set] at (0,-2) {$\{b\}$};
\node[set] at (2,-2) {$\{c\}$};

% 層2
\node[level] (L2) at (0,-4) {};
\node[left] at (-4,-4) {層 2:};
\node[set] at (-3,-4) {$\{a,b\}$};
\node[set] at (-1,-4) {$\{a,c\}$};
\node[set] at (1,-4) {$\{b,c\}$};
\node[set] at (3,-4) {\small $|s| \leq 2$};

% 層3
\node[level] (L3) at (0,-6) {};
\node[left] at (-4,-6) {層 3:};
\node[set] at (0,-6) {$\{a,b,c\}$};
\node[set] at (2.5,-6) {\small $|s| \leq 3$};

\draw[->,thick] (L0) -- (L1);
\draw[->,thick] (L1) -- (L2);
\draw[->,thick] (L2) -- (L3);
\end{tikzpicture}
\end{center}

\subsection{重要な定理}

\begin{proposition}[和集合の濃度]
有限集合 $s, t$ に対して、
\[
\text{minLayer}(s \cup t) \leq \text{minLayer}(s) + \text{minLayer}(t)
\]
すなわち、$|s \cup t| \leq |s| + |t|$ が成り立つ。
\end{proposition}

この定理は、構造塔の枠組みで濃度の基本性質を表現している。

\newpage

\section{具体例3：部分群の包含関係による構造塔}

\subsection{群論の復習}

群 $(G, \cdot, e)$ に対して、部分群 $H \subseteq G$ は以下を満たす：
\begin{enumerate}
\item $e \in H$（単位元を含む）
\item $h_1, h_2 \in H \Rightarrow h_1 \cdot h_2 \in H$（積で閉じている）
\item $h \in H \Rightarrow h^{-1} \in H$（逆元で閉じている）
\end{enumerate}

\subsection{構成}

群 $G$ に対して：

\begin{itemize}
\item $X = G$（群の元）
\item $I = \mathsf{Subgroup}(G)$（部分群全体、包含関係で半順序）
\item $\text{layer}(H) = \{g \in G : g \in H\}$（部分群 $H$ の元）
\item $\text{minLayer}(g) = \langle g \rangle$（$g$ で生成される巡回部分群）
\end{itemize}

\subsection{視覚化（整数加法群の例）}

\begin{center}
\begin{tikzcd}[row sep=large, column sep=large]
& \mathbb{Z} \arrow[d, "\supseteq"] & \text{（全体群）} \\
& 2\mathbb{Z} \arrow[d, "\supseteq"] & \text{（偶数）} \\
& 4\mathbb{Z} \arrow[d, "\supseteq"] & \text{（4の倍数）} \\
& 8\mathbb{Z} & \text{（8の倍数）}
\end{tikzcd}
\end{center}

\begin{examplebox}
整数加法群 $(\mathbb{Z}, +, 0)$ において：
\begin{itemize}
\item $\text{minLayer}(3) = \langle 3 \rangle = 3\mathbb{Z} = \{\ldots, -6, -3, 0, 3, 6, \ldots\}$
\item $\text{minLayer}(6) = \langle 6 \rangle = 6\mathbb{Z} \subseteq 3\mathbb{Z}$
\item $6 \in 3\mathbb{Z}$ だが、$3 \notin 6\mathbb{Z}$
\end{itemize}
\end{examplebox}

\subsection{Lean 4 での実装}

\begin{lstlisting}[style=leanstyle, caption=部分群構造塔の定義]
def SubgroupTower (G : Type*) [Group G] :
    StructureTowerMin G (Subgroup G) where
  layer H := {g : G | g ∈ H}
  covering := by
    intro g
    use ⊤  -- 全体群
    trivial
  monotone := by
    intro H K hHK g hg
    exact hHK hg
  minLayer := fun g => Subgroup.closure {g}
  minLayer_mem := by
    intro g
    exact Subgroup.subset_closure (by simp)
  minLayer_minimal := by
    intro g H hg
    refine (Subgroup.closure_le (K := H)).2 ?_
    intro x hx
    have hx' : x = g := by simpa [Set.mem_singleton_iff] using hx
    simpa [hx'] using hg
\end{lstlisting}

\newpage

\section{具体例4：2の冪による構造塔}

\subsection{構成}

自然数 $\mathbb{N}$ に対して：

\begin{itemize}
\item $X = \mathbb{N}$
\item $I = \mathbb{N}$
\item $\text{layer}(n) = \{k \in \mathbb{N} : k \leq 2^n\}$
\item $\text{minLayer}(k) = \min\{n \in \mathbb{N} : k \leq 2^n\}$
\end{itemize}

\subsection{重要な補題}

\begin{lemma}[被覆性の証明]
任意の $x \in \mathbb{N}$ に対して、$x \leq 2^x$ が成り立つ。
\end{lemma}

\begin{proof}
帰納法で示す。
\begin{itemize}
\item $x = 0$ のとき：$0 \leq 2^0 = 1$ より明らか。
\item $x = k$ で成立すると仮定：$k \leq 2^k$
\item $x = k+1$ のとき：
\[
k + 1 \leq 2^k + 1 \leq 2^k + 2^k = 2 \cdot 2^k = 2^{k+1}
\]
\end{itemize}
よって、すべての $x$ に対して $x \leq 2^x$ が成り立つ。
\end{proof}

\subsection{視覚化}

\begin{center}
\begin{tikzpicture}[scale=0.7]
% 層の描画
\foreach \n/\y in {0/0, 1/1, 2/2, 3/3, 4/4} {
  \pgfmathsetmacro{\pow}{int(2^\n)}
  \draw[fill=maincolor!\n0!white] (-0.5,\y-0.3) rectangle (\pow*0.3+0.5,\y+0.3);
  \node[left] at (-0.5,\y) {層 \n: $[0, 2^\n] = [0, \pow]$};

  % 点の配置
  \foreach \k in {0,...,\pow} {
    \fill (\k*0.3,\y) circle (2pt);
  }
}

% 注釈
\draw[->,thick] (5,0.5) -- (5,3.5) node[midway,right,text width=3cm] {
  指数関数的に\\
  層が拡大
};
\end{tikzpicture}
\end{center}

\begin{examplebox}
具体的な計算：
\begin{itemize}
\item $\text{minLayer}(5) = 3$（なぜなら $5 \leq 2^3 = 8$ だが $5 > 2^2 = 4$）
\item $\text{minLayer}(8) = 3$（$8 = 2^3$）
\item $\text{minLayer}(9) = 4$（$9 \leq 2^4 = 16$ だが $9 > 2^3 = 8$）
\end{itemize}
\end{examplebox}

\newpage

\section{Bourbaki的母構造との対応}

\subsection{3つの母構造における構造塔}

前節までの具体例は、Bourbakiの3つの母構造すべてに構造塔が適用可能であることを示している：

\begin{center}
\begin{tabular}{|c|c|c|c|}
\hline
\textbf{母構造} & \textbf{具体例} & \textbf{添字集合 $I$} & \textbf{層の定義} \\
\hline
代数構造 & 部分群塔 & $\mathsf{Subgroup}(G)$ & 部分群の元 \\
& 主イデアル塔 & $\mathsf{Ideal}(R)$ & イデアルの元 \\
\hline
順序構造 & 整数絶対値塔 & $\mathbb{N}$ & 絶対値以下の整数 \\
& 2の冪塔 & $\mathbb{N}$ & $2^n$ 以下の自然数 \\
\hline
位相構造 & 離散位相塔 & $\mathcal{P}(X)$ & 有限集合 \\
& 位相細分化塔 & $\mathsf{Topology}(X)$ & 開集合 \\
\hline
\end{tabular}
\end{center}

\subsection{母構造の組み合わせ}

Bourbakiの真髄は、母構造を**組み合わせる**ことにある。構造塔もこの観点を自然に扱える：

\subsubsection{順序位相空間}

位相と順序の両方を持つ空間 $(X, \mathcal{T}, \leq)$ において：
\[
\text{layer}(x) = \overline{\{y \in X : y \leq x\}}
\]
（下集合の閉包）

\begin{center}
\begin{tikzcd}[row sep=large]
\text{順序構造} \arrow[r, "\text{closure}"] & \text{位相構造} \\
\{y : y \leq x\} \arrow[r, mapsto] & \overline{\{y : y \leq x\}}
\end{tikzcd}
\end{center}

\subsubsection{位相群}

位相と群構造を持つ空間 $(G, \mathcal{T}, \cdot)$ において：
\[
\text{layer}(U) = \langle U \rangle_{\text{subgroup}}
\]
（開集合 $U$ で生成される部分群）

\begin{remarkbox}
これらの組み合わせは、Lie群論、代数的位相幾何学など、現代数学の基礎を形成している。構造塔は、これらを統一的に扱う枠組みを提供する。
\end{remarkbox}

\newpage

\section{圏論的視点：構造塔の射}

\subsection{構造塔間の射}

2つの構造塔 $(X_1, I_1, \text{layer}_1)$ と $(X_2, I_2, \text{layer}_2)$ の間の**射**（morphism）は、以下を満たす写像 $f : X_1 \to X_2$ である：

\begin{enumerate}
\item $\forall x \in X_1, \ \text{minLayer}_2(f(x)) \leq \phi(\text{minLayer}_1(x))$
\end{enumerate}

ここで、$\phi : I_1 \to I_2$ は順序を保つ写像。

\subsection{具体例：絶対値から濃度への写像}

整数絶対値塔から有限集合濃度塔への射を考える：

\begin{center}
\begin{tikzcd}[row sep=large, column sep=huge]
\mathbb{Z} \arrow[r, "f"] \arrow[d, "\text{minLayer}_1"'] & \mathsf{Finset}(\mathbb{N}) \arrow[d, "\text{minLayer}_2"] \\
\mathbb{N} \arrow[r, "\phi"'] & \mathbb{N}
\end{tikzcd}
\end{center}

例：$f(k) = \{0, 1, \ldots, |k|\}$ とすると、
\[
\text{minLayer}_2(f(k)) = |k| + 1 \quad \text{vs} \quad \text{minLayer}_1(k) = |k|
\]

\subsection{関手性}

構造塔全体は**圏**（category）を形成し、射は関手的性質を持つ：
\begin{itemize}
\item 恒等射 $\text{id}_X : X \to X$ は構造塔の射
\item 射の合成は結合的
\end{itemize}

これにより、構造塔は現代数学の言語（圏論）と自然に接続される。

\newpage

\section{応用と展望}

\subsection{測度論への応用}

測度空間 $(X, \Sigma, \mu)$ において、$\sigma$-代数の構造塔：
\begin{itemize}
\item $I = \{\text{$\sigma$-代数} \subseteq \mathcal{P}(X)\}$
\item $\text{layer}(\Sigma) = \{A \in X : A \in \Sigma\}$（可測集合）
\item $\text{minLayer}(A) = \sigma(A)$（$A$ で生成される最小 $\sigma$-代数）
\end{itemize}

\subsection{確率論への応用（フィルトレーション）}

確率空間 $(\Omega, \mathcal{F}, P)$ 上のフィルトレーション $\{\mathcal{F}_t\}_{t \geq 0}$ は、構造塔の特殊な場合：
\begin{itemize}
\item $I = \mathbb{R}_{\geq 0}$（時間）
\item $\text{layer}(t) = \mathcal{F}_t$（時刻 $t$ までの情報）
\item $s \leq t \Rightarrow \mathcal{F}_s \subseteq \mathcal{F}_t$（単調性）
\end{itemize}

\begin{remarkbox}[title=最新の研究との接続]
本プロジェクトの開発者は、確率論におけるフィルトレーション理論の形式化も進めており、構造塔との統合が期待されている。
\end{remarkbox}

\subsection{代数幾何への応用}

層（sheaf）理論における構造塔：
\begin{itemize}
\item $I = \mathsf{Open}(X)$（開集合、包含関係の逆順序）
\item $\text{layer}(U) = \mathcal{F}(U)$（開集合 $U$ 上の切断）
\item 制限写像による単調性
\end{itemize}

\newpage

\section{Lean 4 による形式化の意義}

\subsection{なぜ形式化するのか}

数学の形式化には以下の利点がある：

\begin{enumerate}
\item \textbf{厳密性の保証}：すべての証明が機械的に検証可能
\item \textbf{一般化の発見}：抽象的なパターンを明示化
\item \textbf{教育的価値}：段階的な理解を促進
\item \textbf{再利用性}：証明を積み重ねて複雑な定理を構築
\end{enumerate}

\subsection{Lean 4 の特徴}

Lean 4 は以下の特徴を持つ定理証明支援系：
\begin{itemize}
\item \textbf{型理論ベース}：Calculus of Inductive Constructions（CIC）
\item \textbf{強力なライブラリ}：Mathlib（数学の標準ライブラリ）
\item \textbf{タクティクシステム}：証明の自動化をサポート
\item \textbf{VSCode統合}：開発環境での即時フィードバック
\end{itemize}

\subsection{本プロジェクトの成果}

\begin{center}
\begin{tabular}{|l|c|c|}
\hline
\textbf{ファイル名} & \textbf{内容} & \textbf{状態} \\
\hline
\texttt{CAT2\_complete.lean} & 構造塔の抽象理論 & ✅ 完成 \\
\texttt{CompleteExamples.lean} & 5つの完全実装例 & ✅ 完成 \\
\texttt{ConcreteExamples.lean} & Bourbaki的母構造 & ⚠️ 一部 sorry \\
\hline
\end{tabular}
\end{center}

\textbf{完全に形式化された構造塔}：
\begin{enumerate}
\item \texttt{IntAbsoluteTower}：整数の絶対値階層
\item \texttt{FinsetCardinalityTower}：有限集合の濃度
\item \texttt{SubgroupTower}：部分群の包含関係
\item \texttt{PowerOfTwoTower}：2の冪による階層
\item \texttt{ListLengthTower}：リストの長さ
\end{enumerate}

\newpage

\section{演習問題}

\subsection{基礎問題}

\begin{enumerate}
\item 整数絶対値塔において、$\text{layer}(5)$ に属する整数をすべて列挙せよ。

\item 有限集合濃度塔において、$\{1, 2\}$ と $\{2, 3\}$ の最小層を求めよ。

\item 部分群塔（整数加法群）において、$\text{minLayer}(12)$ を求め、それに属するすべての整数を述べよ。
\end{enumerate}

\subsection{発展問題}

\begin{enumerate}
\setcounter{enumi}{3}
\item 2の冪塔において、$\text{minLayer}(100)$ を計算せよ。また、その値の意味を説明せよ。

\item 構造塔の射 $f : \text{IntAbsoluteTower} \to \text{FinsetCardinalityTower}$ を1つ構成せよ。

\item 測度論における $\sigma$-代数の構造塔を形式的に定義し、最小層の存在を議論せよ。
\end{enumerate}

\subsection{研究課題}

\begin{enumerate}
\setcounter{enumi}{6}
\item 位相空間における開集合の構造塔を定義し、連続写像との関係を調べよ。

\item フィルトレーションを構造塔として形式化し、マルチンゲール理論への応用を考察せよ。

\item Lean 4 で新しい構造塔の例を実装し、その性質を証明せよ。
\end{enumerate}

\newpage

\section{まとめ}

\subsection{本稿で学んだこと}

\begin{itemize}
\item \textbf{構造塔の定義}：階層的包含関係を統一的に扱う枠組み
\item \textbf{具体例}：整数、有限集合、部分群、2の冪など
\item \textbf{Bourbaki的母構造との対応}：代数、順序、位相すべてに適用可能
\item \textbf{圏論的視点}：射と関手による抽象化
\item \textbf{応用}：測度論、確率論、代数幾何への拡張
\item \textbf{形式化の意義}：Lean 4 による厳密な証明
\end{itemize}

\subsection{構造塔の本質}

構造塔は、以下の3つの要素を統合する：

\begin{center}
\begin{tikzpicture}[
  concept/.style={rectangle, draw=maincolor, fill=maincolor!20, rounded corners, minimum width=3cm, minimum height=1cm, align=center}
]
\node[concept] (A) at (0,0) {階層性\\（順序構造）};
\node[concept] (B) at (4,2) {被覆性\\（すべてを分類）};
\node[concept] (C) at (4,-2) {最小性\\（一意な分類）};

\draw[->,thick] (A) -- (B) node[midway,above,sloped] {単調性};
\draw[->,thick] (A) -- (C) node[midway,below,sloped] {最小層};
\draw[->,thick] (B) -- (C) node[midway,right] {一意性};
\end{tikzpicture}
\end{center}

\subsection{今後の展望}

構造塔理論は、以下の方向に発展する可能性がある：

\begin{enumerate}
\item \textbf{Mathlib への貢献}：Lean の標準ライブラリへの統合
\item \textbf{高階圏論}：$n$-圏における構造塔の一般化
\item \textbf{ホモトピー型理論}：型理論的視点からの再構成
\item \textbf{教育への応用}：数学的構造の統一的教授法
\end{enumerate}

\begin{remarkbox}[title=Bourbakiの夢]
Bourbakiは、数学を「構造の建築学」として再構築することを夢見た。構造塔は、その夢を21世紀のコンピュータサイエンスの力を借りて実現する試みである。すべての数学的対象を、階層的で統一的な言語で語ること――これが構造塔理論の目標である。
\end{remarkbox}

\newpage

\section*{参考文献}

\begin{thebibliography}{99}

\bibitem{bourbaki1968}
N. Bourbaki, \textit{Theory of Sets}, Springer-Verlag, 1968.

\bibitem{mathlib}
The mathlib Community, \textit{The Lean Mathematical Library},
\url{https://github.com/leanprover-community/mathlib4}, 2024.

\bibitem{lean4}
Leonardo de Moura et al., \textit{The Lean 4 Theorem Prover},
\url{https://lean-lang.org/}, 2024.

\bibitem{cat2}
本プロジェクト, \textit{CAT2\_complete.lean: 構造塔の圏論的形式化},
Lean Natural Numbers Project, 2025.

\bibitem{examples}
本プロジェクト, \textit{StructureTower\_CompleteExamples.lean: 完全実装例},
Lean Natural Numbers Project, 2025.

\bibitem{concrete}
本プロジェクト, \textit{StructureTower\_ConcreteExamples.lean: Bourbaki的母構造},
Lean Natural Numbers Project, 2025.

\bibitem{guide}
本プロジェクト, \textit{COMPLETE\_STRENGTHENING\_GUIDE.md: 補強戦略},
Lean Natural Numbers Project, 2025.

\end{thebibliography}

\newpage

\appendix

\section{Lean 4 基本文法ガイド}

\subsection{構造体の定義}

\begin{lstlisting}[style=leanstyle]
structure StructureTowerMin (X : Type*) (I : Type*) [Preorder I] where
  layer : I → Set X
  covering : ∀ x : X, ∃ i : I, x ∈ layer i
  monotone : ∀ {i j : I}, i ≤ j → layer i ⊆ layer j
\end{lstlisting}

\subsection{証明の基本タクティク}

\begin{itemize}
\item \texttt{intro}：仮定を導入
\item \texttt{exact}：完全に一致する項を与える
\item \texttt{simp}：簡約化
\item \texttt{refine}：部分的な項を与え、残りをゴールとする
\item \texttt{by}：タクティクモードに入る
\end{itemize}

\section{用語集}

\begin{description}
\item[構造塔 (Structure Tower)] 階層的包含関係を抽象化した数学的構造
\item[層 (Layer)] 構造塔における各階層の集合
\item[最小層 (Minimal Layer)] 要素が属する最小の層を与える関数
\item[母構造 (Mother Structure)] Bourbakiが提唱した基本的な数学的構造（代数・順序・位相）
\item[形式化 (Formalization)] 数学的議論をコンピュータで検証可能な形式に変換すること
\item[定理証明支援系 (Proof Assistant)] 証明を支援・検証するソフトウェア（Lean, Coq など）
\end{description}

\vfill

\begin{center}
\rule{0.8\textwidth}{0.4pt}\\[1em]
\textit{このドキュメントは、Claude Code (Anthropic) により生成されました。}\\
\textit{Lean 4 形式化は CODEX (OpenAI) により実装されました。}\\
\textit{Lean Natural Numbers - Structure Tower Formalization Project}\\
\textit{2025年}
\end{center}

\end{document}
